\subsubsection*{Overview}

Over the past decade, researchers have been able to formally
verify various key components of computer systems, including compilers,
OS kernels, file systems, and processor designs.
Building on these
successes, the research community is attempting to construct
large-scale certified systems by using formal {\em deep}
specifications as interfaces between the correctness proofs of various
components.
Because these components are different in nature,
the semantic models and proof techniques
used to verify them
are different as well,
complicating their reuse for the construction of
end-to-end correctness proofs.

While researchers have been able to interface
differing semantic models on a case-by-case basis,
and have developed approaches for combining
models with similar foundations
(e.g., multi-language techniques),
these efforts remain ad-hoc or limited in scope.
For this reason,
the construction of end-to-end proofs of correctness
for heterogeneous systems has been restricted to
heroic one-off undertakings.

To overcome these limitations,
we seek to develop new foundations for
coarse-grained, large-scale composition of
certified software components.
Our proposed methodology,
tentatively named \emph{algebraic game semantics},
integrates aspects of
algebraic effects, universal coalgebra and game semantics.
Rather than putting forward
a single semantic framework with a large feature set,
our goal is to allow \emph{existing semantic spaces}
to be described within a unified paradigm;
to allow new ones to be derived from desired features;
and for these various semantic spaces to easily embed into one another.
This approach will facilitate interoperability
between proof techniques and semantic models used by
a variety of software verification developments and tools.

\vspace{+2mm}
\noindent{\bf Keywords}:~~{%
Compositional Semantics;
Algebraic Effects;
Game Semantics;
Compositional Program Verification;
Verified C Compiler;
Deep Specification.}

\subsubsection*{Intellectual Merit}

The project will make three related scientific contributions.
First,
we will develop the theoretical foundations of algebraic game semantics.
These foundations will support
various combinations of features found in existing models, including:
non-termination, non-determinism, refinement,
concurrency, probabilistic behavior, certified compilation, and
information flow security.
Second,
we will mechanize these foundations
as a library for the Rocq prover,
enabling their use for concrete, large-scale verification tasks.
We will use our methodology to
characterize existing Rocq-based semantic frameworks---%
including Interaction Trees and CompCert transition systems---%
making them easy to interface with one another.
Third,
to evaluate the capabilities of our approach,
we will use our library to model a Unix-like environment,
providing semantics and verification capabilities
for process interaction and shell scripting.
Under a gradual verification approach,
we will provide certified versions of
several commands available in this environment,
implemented and verified in C,
and compiled with CompCert.

\subsubsection*{Broader Impacts}

The proposed project aims to develop
new methodologies and tools
enabling the reuse of certified software components
for engineering certified heterogeneous systems.
These new technologies
will greatly facilitate the verification
of large-scale software systems, which in turn will have a profound impact
on the software industry and the society in general. The applicability
of the research outcome can be easily extended to relevant fields such
as operating systems, blockchain and smart contracts, and
cyber-physical systems.  On the educational side, this project will
push new courses on formal semantics, compilers and interpreters, and
language-based security, and will broaden the participation of
underrepresented groups.  Artifacts resulting from the project will be
made open-source to ensure rapid dissemination of ideas.

